% Options for packages loaded elsewhere
% Options for packages loaded elsewhere
\PassOptionsToPackage{unicode}{hyperref}
\PassOptionsToPackage{hyphens}{url}
\PassOptionsToPackage{dvipsnames,svgnames,x11names}{xcolor}
%
\documentclass[
  a4paper,
]{article}
\usepackage{xcolor}
\usepackage[margin=25mm]{geometry}
\usepackage{amsmath,amssymb}
\setcounter{secnumdepth}{5}
\usepackage{iftex}
\ifPDFTeX
  \usepackage[T1]{fontenc}
  \usepackage[utf8]{inputenc}
  \usepackage{textcomp} % provide euro and other symbols
\else % if luatex or xetex
  \usepackage{unicode-math} % this also loads fontspec
  \defaultfontfeatures{Scale=MatchLowercase}
  \defaultfontfeatures[\rmfamily]{Ligatures=TeX,Scale=1}
\fi
\usepackage{lmodern}
\ifPDFTeX\else
  % xetex/luatex font selection
\fi
% Use upquote if available, for straight quotes in verbatim environments
\IfFileExists{upquote.sty}{\usepackage{upquote}}{}
\IfFileExists{microtype.sty}{% use microtype if available
  \usepackage[]{microtype}
  \UseMicrotypeSet[protrusion]{basicmath} % disable protrusion for tt fonts
}{}
\makeatletter
\@ifundefined{KOMAClassName}{% if non-KOMA class
  \IfFileExists{parskip.sty}{%
    \usepackage{parskip}
  }{% else
    \setlength{\parindent}{0pt}
    \setlength{\parskip}{6pt plus 2pt minus 1pt}}
}{% if KOMA class
  \KOMAoptions{parskip=half}}
\makeatother
% Make \paragraph and \subparagraph free-standing
\makeatletter
\ifx\paragraph\undefined\else
  \let\oldparagraph\paragraph
  \renewcommand{\paragraph}{
    \@ifstar
      \xxxParagraphStar
      \xxxParagraphNoStar
  }
  \newcommand{\xxxParagraphStar}[1]{\oldparagraph*{#1}\mbox{}}
  \newcommand{\xxxParagraphNoStar}[1]{\oldparagraph{#1}\mbox{}}
\fi
\ifx\subparagraph\undefined\else
  \let\oldsubparagraph\subparagraph
  \renewcommand{\subparagraph}{
    \@ifstar
      \xxxSubParagraphStar
      \xxxSubParagraphNoStar
  }
  \newcommand{\xxxSubParagraphStar}[1]{\oldsubparagraph*{#1}\mbox{}}
  \newcommand{\xxxSubParagraphNoStar}[1]{\oldsubparagraph{#1}\mbox{}}
\fi
\makeatother


\usepackage{longtable,booktabs,array}
\newcounter{none} % for unnumbered tables
\usepackage{calc} % for calculating minipage widths
% Correct order of tables after \paragraph or \subparagraph
\usepackage{etoolbox}
\makeatletter
\patchcmd\longtable{\par}{\if@noskipsec\mbox{}\fi\par}{}{}
\makeatother
% Allow footnotes in longtable head/foot
\IfFileExists{footnotehyper.sty}{\usepackage{footnotehyper}}{\usepackage{footnote}}
\makesavenoteenv{longtable}
\usepackage{graphicx}
\makeatletter
\newsavebox\pandoc@box
\newcommand*\pandocbounded[1]{% scales image to fit in text height/width
  \sbox\pandoc@box{#1}%
  \Gscale@div\@tempa{\textheight}{\dimexpr\ht\pandoc@box+\dp\pandoc@box\relax}%
  \Gscale@div\@tempb{\linewidth}{\wd\pandoc@box}%
  \ifdim\@tempb\p@<\@tempa\p@\let\@tempa\@tempb\fi% select the smaller of both
  \ifdim\@tempa\p@<\p@\scalebox{\@tempa}{\usebox\pandoc@box}%
  \else\usebox{\pandoc@box}%
  \fi%
}
% Set default figure placement to htbp
\def\fps@figure{htbp}
\makeatother


% definitions for citeproc citations
\NewDocumentCommand\citeproctext{}{}
\NewDocumentCommand\citeproc{mm}{%
  \begingroup\def\citeproctext{#2}\cite{#1}\endgroup}
\makeatletter
 % allow citations to break across lines
 \let\@cite@ofmt\@firstofone
 % avoid brackets around text for \cite:
 \def\@biblabel#1{}
 \def\@cite#1#2{{#1\if@tempswa , #2\fi}}
\makeatother
\newlength{\cslhangindent}
\setlength{\cslhangindent}{1.5em}
\newlength{\csllabelwidth}
\setlength{\csllabelwidth}{3em}
\newenvironment{CSLReferences}[2] % #1 hanging-indent, #2 entry-spacing
 {\begin{list}{}{%
  \setlength{\itemindent}{0pt}
  \setlength{\leftmargin}{0pt}
  \setlength{\parsep}{0pt}
  % turn on hanging indent if param 1 is 1
  \ifodd #1
   \setlength{\leftmargin}{\cslhangindent}
   \setlength{\itemindent}{-1\cslhangindent}
  \fi
  % set entry spacing
  \setlength{\itemsep}{#2\baselineskip}}}
 {\end{list}}
\usepackage{calc}
\newcommand{\CSLBlock}[1]{\hfill\break\parbox[t]{\linewidth}{\strut\ignorespaces#1\strut}}
\newcommand{\CSLLeftMargin}[1]{\parbox[t]{\csllabelwidth}{\strut#1\strut}}
\newcommand{\CSLRightInline}[1]{\parbox[t]{\linewidth - \csllabelwidth}{\strut#1\strut}}
\newcommand{\CSLIndent}[1]{\hspace{\cslhangindent}#1}



\setlength{\emergencystretch}{3em} % prevent overfull lines

\providecommand{\tightlist}{%
  \setlength{\itemsep}{0pt}\setlength{\parskip}{0pt}}



 


\makeatletter
\@ifpackageloaded{caption}{}{\usepackage{caption}}
\AtBeginDocument{%
\ifdefined\contentsname
  \renewcommand*\contentsname{Table of contents}
\else
  \newcommand\contentsname{Table of contents}
\fi
\ifdefined\listfigurename
  \renewcommand*\listfigurename{List of Figures}
\else
  \newcommand\listfigurename{List of Figures}
\fi
\ifdefined\listtablename
  \renewcommand*\listtablename{List of Tables}
\else
  \newcommand\listtablename{List of Tables}
\fi
\ifdefined\figurename
  \renewcommand*\figurename{Figure}
\else
  \newcommand\figurename{Figure}
\fi
\ifdefined\tablename
  \renewcommand*\tablename{Table}
\else
  \newcommand\tablename{Table}
\fi
}
\@ifpackageloaded{float}{}{\usepackage{float}}
\floatstyle{ruled}
\@ifundefined{c@chapter}{\newfloat{codelisting}{h}{lop}}{\newfloat{codelisting}{h}{lop}[chapter]}
\floatname{codelisting}{Listing}
\newcommand*\listoflistings{\listof{codelisting}{List of Listings}}
\makeatother
\makeatletter
\makeatother
\makeatletter
\@ifpackageloaded{caption}{}{\usepackage{caption}}
\@ifpackageloaded{subcaption}{}{\usepackage{subcaption}}
\makeatother
\usepackage{bookmark}
\IfFileExists{xurl.sty}{\usepackage{xurl}}{} % add URL line breaks if available
\urlstyle{same}
% fallback for those not using the hyperref driver hyperxmp:
\makeatletter
\@ifundefined{xmpquote}{\newcommand{\xmpquote}[1]{#1}}{}
\makeatother
\hypersetup{
  pdftitle={Tap water alkalinity and coffee acidity: a population-weighted atlas and a prediction model for filter brewing},
  pdfauthor={Fellipe Giuliano Duarte},
  pdfkeywords={\xmpquote{coffee}, \xmpquote{brewing
water}, \xmpquote{alkalinity}, \xmpquote{hardness}, \xmpquote{titratable
acidity}, \xmpquote{open data}},
  colorlinks=true,
  linkcolor={blue},
  filecolor={Maroon},
  citecolor={Blue},
  urlcolor={Blue},
  pdfcreator={LaTeX via pandoc}}


\title{Tap water alkalinity and coffee acidity: a population-weighted
atlas and a prediction model for filter brewing}
\author{Fellipe Giuliano Duarte}
\date{2026-09-14}
\begin{document}
\maketitle
\begin{abstract}
Bicarbonate in brewing water neutralises part of the acidity of coffee.
We quantify this effect for municipal tap water. We compiled an open,
provenance-tracked atlas of finished tap-water chemistry for about
47,000 localities in the United States, England, Northern Ireland and
Brazil, added the label compositions of 187 bottled waters sold in seven
countries, and weighted all statistics by population served. Combining
the atlas with published drip-coffee measurements of titratable acidity
by roast level and of sensory intensity by brew strength and extraction
yield, we predict the share of a coffee's acidity that a given water
neutralises. On a population basis, the median United States tap water
neutralises 13 percent of a light roast's acidity and 16 percent of a
dark roast's; the upper quartile neutralises 21 to 26 percent; about 10
percent of the population receives water above 150 mg/L alkalinity as
CaCO\textsubscript{3}, where all roast levels are predicted to need
compensation. Alkalinity and hardness are correlated along one line
across countries, so hardness, which is reported far more widely,
predicts alkalinity within a factor of 1.7. Each 50 mg/L of alkalinity
above the Specialty Coffee Association target neutralises about 1 meq/L
of acid and lowers predicted sourness by 3.6 points on a 100-point
scale; below about 100 mg/L this can be offset by 5 to 10 percent more
coffee per 40 mg/L, above 150 mg/L dilution with low-alkalinity water
requires less change than any recipe adjustment. Whether the shift is
unwelcome depends on the drinker: most black-coffee consumers in
published preference studies dislike sourness. The predictions carry
stated uncertainties and have not been tested by tasting; a paired-brew
protocol is provided. Data, code and figures are open.
\end{abstract}


\section{Introduction}\label{introduction}

Municipal tap water varies widely in hardness (calcium and magnesium),
alkalinity (mostly bicarbonate) and pH. Published chemistry describes
how these variables affect coffee extraction and flavour. Bicarbonate
neutralises the organic and chlorogenic acids that give coffee its
perceived acidity (Yeager et al., 2023); dissolved cations may change
the extraction of some compounds (Bratthäll et al., 2024; Hendon et al.,
2014). These two bodies of knowledge have not been combined at scale: a
map of the water people brew with, the chemistry applied to it, and a
statement for each place and roast level of what the water does to the
cup and how a brewer can compensate.

This study reports model predictions, not new sensory measurements. The
predictions are stated in a form that can be tested, and Section 6
describes the test.

We address three questions.

\begin{enumerate}
\def\labelenumi{\arabic{enumi}.}
\tightlist
\item
  How are alkalinity and hardness distributed in tap water when weighted
  by the population served, and can alkalinity be predicted where only
  hardness is reported?
\item
  What share of a coffee's acidity does a given water neutralise, by
  roast level, and for what share of the population is this share large?
\item
  What does compensation cost, in coffee dose, in bitterness, or in
  blended water?
\end{enumerate}

\section{Data}\label{data}

\subsection{Sources}\label{sources}

Table~\ref{tbl-sources} lists the sources in the first version of the
atlas. Each row of the underlying measurement table records the source
file, its checksum, the download date, the licence, the original
parameter name, value and unit, and the conversion factor applied.

\begin{longtable}[]{@{}
  >{\raggedright\arraybackslash}p{(\linewidth - 8\tabcolsep) * \real{0.2000}}
  >{\raggedright\arraybackslash}p{(\linewidth - 8\tabcolsep) * \real{0.2000}}
  >{\raggedright\arraybackslash}p{(\linewidth - 8\tabcolsep) * \real{0.2000}}
  >{\raggedright\arraybackslash}p{(\linewidth - 8\tabcolsep) * \real{0.2000}}
  >{\raggedright\arraybackslash}p{(\linewidth - 8\tabcolsep) * \real{0.2000}}@{}}
\caption{Sources in atlas v0. A locality is a city, utility, small area,
municipality or bottled product; counts are after aggregation to one row
per locality.}\label{tbl-sources}\tabularnewline
\toprule\noalign{}
\begin{minipage}[b]{\linewidth}\raggedright
source
\end{minipage} & \begin{minipage}[b]{\linewidth}\raggedright
localities
\end{minipage} & \begin{minipage}[b]{\linewidth}\raggedright
with hardness
\end{minipage} & \begin{minipage}[b]{\linewidth}\raggedright
with alkalinity
\end{minipage} & \begin{minipage}[b]{\linewidth}\raggedright
licence
\end{minipage} \\
\midrule\noalign{}
\endfirsthead
\toprule\noalign{}
\begin{minipage}[b]{\linewidth}\raggedright
source
\end{minipage} & \begin{minipage}[b]{\linewidth}\raggedright
localities
\end{minipage} & \begin{minipage}[b]{\linewidth}\raggedright
with hardness
\end{minipage} & \begin{minipage}[b]{\linewidth}\raggedright
with alkalinity
\end{minipage} & \begin{minipage}[b]{\linewidth}\raggedright
licence
\end{minipage} \\
\midrule\noalign{}
\endhead
\bottomrule\noalign{}
\endlastfoot
Brazil, SISAGUA semestral control (Ministério da Saúde) & 2,473 & 2,441
& 0 & CC BY \\
United States, TapWaterData compiled city hardness & 7,614 & 7,614 & 0 &
CC BY 4.0 \\
United States, EPA Six-Year Review 4 finished-water monitoring 2012 to
2019 & 27,030 & 0 & 15,240 & public domain \\
United States, the two above joined on utility id & 4,893 & 4,893 &
3,341 & derived \\
England, water-industry open-data hub (Wessex, Southern, Yorkshire) &
5,580 & 3,839 & 2,396 & CC BY 4.0 \\
Northern Ireland Water customer-tap results & 64 & 64 & 0 & OGL v3 \\
Bottled waters, label compositions (7 countries) & 171 & 146 & 142 & CC
BY 4.0 compilation \\
\end{longtable}

\textbf{Brazil.} The national drinking-water surveillance system SISAGUA
publishes the semestral control samples that suppliers report
(Ministério da Saúde do Brasil, 2026). Total hardness, sodium and pH are
present; alkalinity, calcium and magnesium are not measured in the
system. We loaded the files for 2023 to 2026 (8.2 million rows) and kept
194,226 measurements of the three variables at treatment outlets,
distribution systems and consumption points in 2,441 municipalities.
Source type (surface or ground) comes from the intake registry.

\textbf{United States.} TapWaterData compiles utility-reported hardness
for 16,561 cities from Consumer Confidence Reports, with a tier flag
that separates utility-measured values from county-level ambient
estimates (TapWaterData, 2026); we use the measured tiers only (7,614
cities). The EPA Six-Year Review 4 compliance data give finished-water
alkalinity and pH per public water system for 2012 to 2019 (United
States Environmental Protection Agency, 2024). Joining the two datasets
on the utility identifier gives hardness and alkalinity for 3,341 cities
served by 1,670 utilities. The Water Quality Portal was examined and not
used: its 1.3 million recent results are ambient surface water and
groundwater, with fewer than a dozen finished-water sites.

\textbf{England and Northern Ireland.} The water industry's open-data
portal publishes tap-sample results per small area (Lower-layer Super
Output Area, about 1,500 residents) for every company in England and
Wales (Stream, 2026). Three companies report the variables needed here:
Wessex Water (alkalinity and hardness, with units), Southern Water
(calcium, magnesium and alkalinity, with an empty units column) and
Yorkshire Water (hardness). Northern Ireland Water publishes hardness,
calcium, magnesium, sodium and pH per supply zone under the Open
Government Licence. The other companies publish alkalinity only in
per-postcode lookups whose terms of use reserve all rights; those values
were not used.

\textbf{Bottled water.} We compiled the label compositions of 187
bottled waters sold in supermarkets in Brazil, the United States, the
United Kingdom, Germany, France, Italy and Portugal from brand websites
and official label images, with one source URL per product; 182 have an
official source. Alkalinity was computed as bicarbonate multiplied by
0.820 where bicarbonate is printed.

\subsection{Units and reporting basis}\label{units-and-reporting-basis}

Hardness and alkalinity are stored in mg/L as CaCO\textsubscript{3},
ions in mg/L of the ion, with conversion factors derived from molar
masses. Three sources report ``total hardness'' in mg/L where the figure
is in fact mg/L as calcium, a factor of 2.5 lower than mg/L as
CaCO\textsubscript{3}: Northern Ireland Water (the CaCO\textsubscript{3}
column of the zone report equals the CSV figure multiplied by 2.4985),
Southern Water (the hardness column equals the calcium column, and the
company's online checker labels its default unit ``Cal mg/l'') and SES
Water. Wherever calcium and magnesium are both reported, hardness is
derived from them and the reported hardness column is not used. Every
value whose unit basis had to be assumed is flagged.

\subsection{Population weighting}\label{population-weighting}

A count of localities overstates the alkalinity of the water most people
receive, because many small groundwater systems each count once. We
attached a population to every locality: people served for United States
utilities, mid-2024 estimates for English small areas, and the 2022
census for Brazilian municipalities. The per-locality median alkalinity
of finished water in the United States and England is 187 mg/L; the
population-weighted median is 83 mg/L. All results below are
population-weighted unless stated otherwise.

\subsection{Alkalinity imputation}\label{alkalinity-imputation}

Brazil has no alkalinity data. Across 5,658 United States and English
localities with both variables, a regression of log alkalinity on log
hardness and source type has R\textsuperscript{2} 0.48 and a residual
standard deviation of 0.52 on the log scale, which corresponds to a
multiplicative uncertainty of 1.7 in either direction. We applied this
regression to Brazil's 2,441 municipalities and flagged every imputed
value. Headline statistics are reported with and without Brazil.

\section{Model}\label{model}

\subsection{Residual acidity}\label{residual-acidity}

Coffee acids protonate the bicarbonate in brewing water. At brew pH 4.9
to 5.1, 94 to 96 percent of bicarbonate is protonated (carbonic acid
pK\textsubscript{a1} 6.35), so the neutralisation is close to one
milliequivalent of acid per milliequivalent of alkalinity. We define the
residual acidity of a brew in meq/L as

\[
\mathrm{residual} = \mathrm{TA}_\mathrm{intrinsic}(\mathrm{roast}, \mathrm{TDS}) - f(\mathrm{pH}) \cdot \frac{\mathrm{alkalinity}}{50.04},
\]

where TA is the titratable acidity to pH 8.2, alkalinity is in mg/L as
CaCO\textsubscript{3}, and \(f\) is the protonated fraction of
bicarbonate at the brew pH. The neutralised share is the second term
divided by the first.

\subsection{Parameters}\label{parameters}

Titratable acidity by roast level is taken from the one drip-coffee
study that measured it across brew strength and extraction yield (Batali
et al., 2021): 258 brews of one washed Honduran Arabica at three roast
levels, three strengths (1.0, 1.25 and 1.5 percent total dissolved
solids) and three extraction yields (16, 20 and 24 percent). Titratable
acidity is linear in strength and nearly independent of extraction
yield. Table~\ref{tbl-params} gives the values. The brew water in that
study contained 32 mg/L alkalinity as CaCO\textsubscript{3} (computed
from the printed mineral recipe), so the measured acidity had already
lost about 0.6 meq/L to bicarbonate; the model adds this amount back
before subtracting the alkalinity of the target water.

\begin{longtable}[]{@{}
  >{\raggedright\arraybackslash}p{(\linewidth - 10\tabcolsep) * \real{0.1667}}
  >{\raggedright\arraybackslash}p{(\linewidth - 10\tabcolsep) * \real{0.1667}}
  >{\raggedright\arraybackslash}p{(\linewidth - 10\tabcolsep) * \real{0.1667}}
  >{\raggedright\arraybackslash}p{(\linewidth - 10\tabcolsep) * \real{0.1667}}
  >{\raggedright\arraybackslash}p{(\linewidth - 10\tabcolsep) * \real{0.1667}}
  >{\raggedright\arraybackslash}p{(\linewidth - 10\tabcolsep) * \real{0.1667}}@{}}
\caption{Model parameters per roast bin, from Batali et al.~(2021) Table
1 and Supplemental Figure 1. The brew water in that study had 32 mg/L
alkalinity as CaCO\textsubscript{3} (computed from the printed
recipe).}\label{tbl-params}\tabularnewline
\toprule\noalign{}
\begin{minipage}[b]{\linewidth}\raggedright
roast
\end{minipage} & \begin{minipage}[b]{\linewidth}\raggedright
titratable acidity, meq/L (measured, TDS 1.25 \%)
\end{minipage} & \begin{minipage}[b]{\linewidth}\raggedright
range
\end{minipage} & \begin{minipage}[b]{\linewidth}\raggedright
intrinsic (brew water alkalinity added back)
\end{minipage} & \begin{minipage}[b]{\linewidth}\raggedright
brew pH
\end{minipage} & \begin{minipage}[b]{\linewidth}\raggedright
protonated share of bicarbonate
\end{minipage} \\
\midrule\noalign{}
\endfirsthead
\toprule\noalign{}
\begin{minipage}[b]{\linewidth}\raggedright
roast
\end{minipage} & \begin{minipage}[b]{\linewidth}\raggedright
titratable acidity, meq/L (measured, TDS 1.25 \%)
\end{minipage} & \begin{minipage}[b]{\linewidth}\raggedright
range
\end{minipage} & \begin{minipage}[b]{\linewidth}\raggedright
intrinsic (brew water alkalinity added back)
\end{minipage} & \begin{minipage}[b]{\linewidth}\raggedright
brew pH
\end{minipage} & \begin{minipage}[b]{\linewidth}\raggedright
protonated share of bicarbonate
\end{minipage} \\
\midrule\noalign{}
\endhead
\bottomrule\noalign{}
\endlastfoot
light & 12.46 & 11.5 to 13.0 & 13.08 & 4.93 & 0.963 \\
medium & 11.36 & 10.4 to 12.0 & 11.98 & 5.01 & 0.956 \\
dark & 9.78 & 9.2 to 10.4 & 10.39 & 5.14 & 0.942 \\
\end{longtable}

\subsection{Cation-dependent
extraction}\label{cation-dependent-extraction}

Hendon and colleagues predicted from binding energies that magnesium and
calcium increase the extraction of coffee acids (Hendon et al., 2014). A
GC-MS and NMR study found no such effect on organic acids at 100 mg/L
(Bratthäll et al., 2024). A three-water filter-brewing study found that
dark roast brewed in the hardest water was rated more acidic by a
recognition panel, which the authors attribute to calcium-assisted
extraction (Liu et al., 2025). None of these studies gives a magnitude
that can be used as a parameter. We include a multiplicative extraction
term on titratable acidity with a central value of zero and an upper
bound of a 5 percent yield increase per meq/L of divalent cations, and
report results for both. The result of Liu and colleagues suggests that
the upper bound is relevant for dark roasts.

\subsection{Alkalinity bands and matching
criterion}\label{alkalinity-bands-and-matching-criterion}

We classify water into five alkalinity bands with edges at 40, 80, 150
and 250 mg/L as CaCO\textsubscript{3}. A Gaussian-mixture clustering was
also fitted; it did not identify separated clusters. Tap water forms a
continuum along the alkalinity--hardness line, and the information
criterion continued to decrease up to the six-component limit. Fixed
bands are equally defensible and simpler to communicate.

For each band and roast we compute the residual acidity at the band's
population-weighted median water relative to the Specialty Coffee
Association reference water (40 mg/L alkalinity, 68 mg/L hardness
(Specialty Coffee Association of America, 2009)). A ratio of 0.75 to
1.25 is classified as \emph{works}, 0.5 to 0.75 as \emph{compensate},
and below 0.5 as \emph{fights}.

\subsection{Sensory model and compensation
rules}\label{sensory-model-and-compensation-rules}

The trained-panel study from which the titratable acidity values come
also reports marginal mean intensities of 32 attributes at each level of
roast, strength and extraction (Frost et al., 2020). The design is
balanced, so an additive model reproduces the marginal means exactly.
Interaction terms are not published, and the same group's response
surface for sourness is a first-order plane (Batali et al., 2021), so
main effects only are used. Sourness increases by 21 points per
percentage point of strength and decreases by 0.83 points per percentage
point of extraction; bitterness increases by 23 points per percentage
point of strength.

Two calibrations connect water chemistry to perception. The roast means
of Batali and colleagues give 3.61 sourness points per meq/L of
titratable acidity. Dividing the acid neutralised by a water by the
slope of titratable acidity against strength gives the strength increase
that restores the acid balance (the chemistry estimate). Dividing the
predicted sourness loss by the slope of sourness against strength gives
the increase that restores the predicted sourness (the sensory
estimate). The two estimates differ by a factor of 1.8, and both are
reported. A third route is dilution: alkalinity is linear in mixing, so
a volume fraction \(1 - 40/\mathrm{alk}\) of zero-alkalinity water
brings a tap water of alkalinity alk to the reference value.

\section{Results}\label{results}

\subsection{Distribution of tap-water alkalinity and
hardness}\label{distribution-of-tap-water-alkalinity-and-hardness}

Figure~\ref{fig-where} shows the 13,933 localities with both variables.
Alkalinity and hardness are correlated along one line in every country.
England's chalk catchments occupy the upper end; the imputed Brazilian
municipalities form straight lines by construction. A small group of
softened waters, with alkalinity retained and hardness removed, lies
above the line and covers about 1 percent of the population.

\begin{figure}[H]

\centering{

\pandocbounded{\includegraphics[keepaspectratio]{figures/fig2_where_waters_sit.png}}

}

\caption{\label{fig-where}Alkalinity against hardness for 13,933
localities. Filled symbols are measured alkalinity; open symbols are
imputed from hardness. Symbol area is proportional to the population
served. Horizontal lines mark the band edges at 40, 80, 150 and 250
mg/L. The star marks the Specialty Coffee Association reference water.}

\end{figure}%

\subsection{Neutralised share of coffee
acidity}\label{neutralised-share-of-coffee-acidity}

Figure~\ref{fig-neutral} shows the model output. At the reference
alkalinity of 40 mg/L the water neutralises 6 percent of a light roast's
titratable acidity and 8 percent of a dark roast's. At the
population-weighted median for the United States (83 mg/L) the shares
are 13 and 16 percent; at the population-weighted upper quartile (136
mg/L) they are 21 and 26 percent. For 20 percent of the population the
water neutralises more than a quarter of a light roast's acidity, and
for 31 percent more than a quarter of a dark roast's. The dark roast
loses the largest share at every alkalinity because it has the lowest
titratable acidity.

\begin{figure}[H]

\centering{

\pandocbounded{\includegraphics[keepaspectratio]{figures/fig1_acid_neutralised.png}}

}

\caption{\label{fig-neutral}Share of the brew's titratable acidity
neutralised by the water, by roast level. The shaded band is the range
of titratable acidity in Batali et al. (2021). The lower panel is the
population-weighted distribution of measured alkalinity in United States
and English tap water.}

\end{figure}%

\subsection{Population by alkalinity
band}\label{population-by-alkalinity-band}

Table~\ref{tbl-bands} and Figure~\ref{fig-bands} give the distribution
by band. Bands one to three (up to 150 mg/L) hold about 90 percent of
the population; the reference alkalinity lies at the upper edge of the
first band. Bands four and five (above 150 mg/L) hold about 10 percent
of the population overall (12 percent when Brazil's imputed values are
excluded), but half of the English population served by the three
companies with data, all of which draw on chalk aquifers.

\begin{longtable}[]{@{}
  >{\raggedright\arraybackslash}p{(\linewidth - 18\tabcolsep) * \real{0.1000}}
  >{\raggedright\arraybackslash}p{(\linewidth - 18\tabcolsep) * \real{0.1000}}
  >{\raggedright\arraybackslash}p{(\linewidth - 18\tabcolsep) * \real{0.1000}}
  >{\raggedright\arraybackslash}p{(\linewidth - 18\tabcolsep) * \real{0.1000}}
  >{\raggedright\arraybackslash}p{(\linewidth - 18\tabcolsep) * \real{0.1000}}
  >{\raggedright\arraybackslash}p{(\linewidth - 18\tabcolsep) * \real{0.1000}}
  >{\raggedright\arraybackslash}p{(\linewidth - 18\tabcolsep) * \real{0.1000}}
  >{\raggedright\arraybackslash}p{(\linewidth - 18\tabcolsep) * \real{0.1000}}
  >{\raggedright\arraybackslash}p{(\linewidth - 18\tabcolsep) * \real{0.1000}}
  >{\raggedright\arraybackslash}p{(\linewidth - 18\tabcolsep) * \real{0.1000}}@{}}
\caption{Alkalinity bands and who brews with them. Population-weighted;
`measured only' excludes Brazil, whose alkalinity is imputed. Country
columns give the share of that country's population in the
band.}\label{tbl-bands}\tabularnewline
\toprule\noalign{}
\begin{minipage}[b]{\linewidth}\raggedright
band (mg/L as CaCO\textsubscript{3})
\end{minipage} & \begin{minipage}[b]{\linewidth}\raggedright
name
\end{minipage} & \begin{minipage}[b]{\linewidth}\raggedright
localities
\end{minipage} & \begin{minipage}[b]{\linewidth}\raggedright
share of people, all
\end{minipage} & \begin{minipage}[b]{\linewidth}\raggedright
share of people, measured only
\end{minipage} & \begin{minipage}[b]{\linewidth}\raggedright
US
\end{minipage} & \begin{minipage}[b]{\linewidth}\raggedright
England
\end{minipage} & \begin{minipage}[b]{\linewidth}\raggedright
Brazil (imputed)
\end{minipage} & \begin{minipage}[b]{\linewidth}\raggedright
median alkalinity
\end{minipage} & \begin{minipage}[b]{\linewidth}\raggedright
median hardness
\end{minipage} \\
\midrule\noalign{}
\endfirsthead
\toprule\noalign{}
\begin{minipage}[b]{\linewidth}\raggedright
band (mg/L as CaCO\textsubscript{3})
\end{minipage} & \begin{minipage}[b]{\linewidth}\raggedright
name
\end{minipage} & \begin{minipage}[b]{\linewidth}\raggedright
localities
\end{minipage} & \begin{minipage}[b]{\linewidth}\raggedright
share of people, all
\end{minipage} & \begin{minipage}[b]{\linewidth}\raggedright
share of people, measured only
\end{minipage} & \begin{minipage}[b]{\linewidth}\raggedright
US
\end{minipage} & \begin{minipage}[b]{\linewidth}\raggedright
England
\end{minipage} & \begin{minipage}[b]{\linewidth}\raggedright
Brazil (imputed)
\end{minipage} & \begin{minipage}[b]{\linewidth}\raggedright
median alkalinity
\end{minipage} & \begin{minipage}[b]{\linewidth}\raggedright
median hardness
\end{minipage} \\
\midrule\noalign{}
\endhead
\bottomrule\noalign{}
\endlastfoot
0 to 40 & very low alkalinity & 2,312 & 20\% & 18\% & 18\% & 1\% & 40\%
& 30 & 34 \\
40 to 80 & low alkalinity & 3,331 & 47\% & 47\% & 48\% & 13\% & 44\% &
64 & 102 \\
80 to 150 & moderate alkalinity & 3,795 & 22\% & 22\% & 23\% & 21\% &
14\% & 108 & 108 \\
150 to 250 & high alkalinity & 3,404 & 10\% & 11\% & 10\% & 50\% & 3\% &
213 & 180 \\
over 250 & very high alkalinity & 1,114 & 1\% & 1\% & 1\% & 15\% & 0\% &
280 & 310 \\
\end{longtable}

\begin{figure}[H]

\centering{

\pandocbounded{\includegraphics[keepaspectratio]{figures/fig3_population_by_band.png}}

}

\caption{\label{fig-bands}Share of each country's population by
alkalinity band. England is represented by the three water companies
with measured alkalinity; Brazilian alkalinity is imputed.}

\end{figure}%

\subsection{Matching table}\label{matching-table}

Table~\ref{tbl-matching} applies the classification criterion. All roast
levels are classified as \emph{works} in bands one to three, at the band
medians and at the band edges. All roast levels are classified as
\emph{compensate} in bands four and five: in band four the ratio ranges
from 0.73 for light roast to 0.66 for dark roast, in band five from 0.62
to 0.53. Differences between roast levels within a band are smaller than
differences between bands.

\begin{longtable}[]{@{}
  >{\raggedright\arraybackslash}p{(\linewidth - 6\tabcolsep) * \real{0.2500}}
  >{\raggedright\arraybackslash}p{(\linewidth - 6\tabcolsep) * \real{0.2500}}
  >{\raggedright\arraybackslash}p{(\linewidth - 6\tabcolsep) * \real{0.2500}}
  >{\raggedright\arraybackslash}p{(\linewidth - 6\tabcolsep) * \real{0.2500}}@{}}
\caption{Matching table v0: residual acidity at each band's median water
relative to the SCA reference water (40 mg/L alkalinity, 68 mg/L
hardness). 0.75 to 1.25 = works; 0.5 to 0.75 = compensate; below 0.5 =
fights.}\label{tbl-matching}\tabularnewline
\toprule\noalign{}
\begin{minipage}[b]{\linewidth}\raggedright
band
\end{minipage} & \begin{minipage}[b]{\linewidth}\raggedright
light roast
\end{minipage} & \begin{minipage}[b]{\linewidth}\raggedright
medium roast
\end{minipage} & \begin{minipage}[b]{\linewidth}\raggedright
dark roast
\end{minipage} \\
\midrule\noalign{}
\endfirsthead
\toprule\noalign{}
\begin{minipage}[b]{\linewidth}\raggedright
band
\end{minipage} & \begin{minipage}[b]{\linewidth}\raggedright
light roast
\end{minipage} & \begin{minipage}[b]{\linewidth}\raggedright
medium roast
\end{minipage} & \begin{minipage}[b]{\linewidth}\raggedright
dark roast
\end{minipage} \\
\midrule\noalign{}
\endhead
\bottomrule\noalign{}
\endlastfoot
1 & 1.02 (works) & 1.02 (works) & 1.02 (works) \\
2 & 0.96 (works) & 0.96 (works) & 0.95 (works) \\
3 & 0.89 (works) & 0.88 (works) & 0.87 (works) \\
4 & 0.73 (compensate (flatter)) & 0.71 (compensate (flatter)) & 0.66
(compensate (flatter)) \\
5 & 0.62 (compensate (flatter)) & 0.59 (compensate (flatter)) & 0.53
(compensate (flatter)) \\
\end{longtable}

\subsection{Compensation rules}\label{compensation-rules}

Figure~\ref{fig-rules} and Table~\ref{tbl-rules} translate the model
into recipe changes. Each 50 mg/L of alkalinity above the reference
neutralises about 1 meq/L of acid and lowers predicted sourness by 3.6
points. Up to about 100 mg/L this can be offset by 5 to 10 percent more
coffee per 40 mg/L above the reference, at a cost of 1 to 2 bitterness
points. At 150 mg/L the strength route requires 16 to 29 percent more
coffee and adds 4.5 bitterness points; at 250 mg/L, 31 to 55 percent and
8.5 points. Restoring sourness by lowering extraction alone would
require a reduction of 9 percentage points in yield at 150 mg/L, whereas
a change of grind size typically moves yield by 2 to 4 points, so
extraction is insufficient on its own. Above 150 mg/L, dilution with
zero-alkalinity water (73 percent at 150 mg/L, 84 percent at 250 mg/L)
or the use of a low-alkalinity bottled water requires less change than
any recipe adjustment.

\begin{longtable}[]{@{}
  >{\raggedright\arraybackslash}p{(\linewidth - 10\tabcolsep) * \real{0.1667}}
  >{\raggedright\arraybackslash}p{(\linewidth - 10\tabcolsep) * \real{0.1667}}
  >{\raggedright\arraybackslash}p{(\linewidth - 10\tabcolsep) * \real{0.1667}}
  >{\raggedright\arraybackslash}p{(\linewidth - 10\tabcolsep) * \real{0.1667}}
  >{\raggedright\arraybackslash}p{(\linewidth - 10\tabcolsep) * \real{0.1667}}
  >{\raggedright\arraybackslash}p{(\linewidth - 10\tabcolsep) * \real{0.1667}}@{}}
\caption{Dial-in rules v0 for a light roast at TDS 1.25 \% and 20 \%
extraction. Medium and dark roasts differ by under 10 \% in every column
(full table in the repository).}\label{tbl-rules}\tabularnewline
\toprule\noalign{}
\begin{minipage}[b]{\linewidth}\raggedright
alkalinity, mg/L
\end{minipage} & \begin{minipage}[b]{\linewidth}\raggedright
acid neutralised vs reference, meq/L
\end{minipage} & \begin{minipage}[b]{\linewidth}\raggedright
predicted sourness change (0 to 100)
\end{minipage} & \begin{minipage}[b]{\linewidth}\raggedright
extra coffee, chemistry to sensory estimate
\end{minipage} & \begin{minipage}[b]{\linewidth}\raggedright
bitterness side effect
\end{minipage} & \begin{minipage}[b]{\linewidth}\raggedright
or blend with zero-alkalinity water
\end{minipage} \\
\midrule\noalign{}
\endfirsthead
\toprule\noalign{}
\begin{minipage}[b]{\linewidth}\raggedright
alkalinity, mg/L
\end{minipage} & \begin{minipage}[b]{\linewidth}\raggedright
acid neutralised vs reference, meq/L
\end{minipage} & \begin{minipage}[b]{\linewidth}\raggedright
predicted sourness change (0 to 100)
\end{minipage} & \begin{minipage}[b]{\linewidth}\raggedright
extra coffee, chemistry to sensory estimate
\end{minipage} & \begin{minipage}[b]{\linewidth}\raggedright
bitterness side effect
\end{minipage} & \begin{minipage}[b]{\linewidth}\raggedright
or blend with zero-alkalinity water
\end{minipage} \\
\midrule\noalign{}
\endhead
\bottomrule\noalign{}
\endlastfoot
20 & -0.39 & +1.4 & -3 to -5 \% & -0.9 & 0 \% \\
40 & +0.00 & -0.0 & +0 to +0 \% & +0.0 & 0 \% \\
80 & +0.77 & -2.8 & +6 to +10 \% & +1.6 & 50 \% \\
120 & +1.54 & -5.6 & +12 to +21 \% & +3.2 & 67 \% \\
150 & +2.12 & -7.6 & +16 to +29 \% & +4.5 & 73 \% \\
200 & +3.08 & -11.1 & +24 to +42 \% & +6.5 & 80 \% \\
250 & +4.04 & -14.6 & +31 to +55 \% & +8.5 & 84 \% \\
\end{longtable}

\begin{figure}[H]

\centering{

\pandocbounded{\includegraphics[keepaspectratio]{figures/fig4_dialin_rules.png}}

}

\caption{\label{fig-rules}Left: predicted change in sourness relative to
the reference water, by roast level. Right: additional coffee required
to compensate; the thick line restores titratable acidity, the thin line
restores predicted sourness. The shaded region marks alkalinities above
150 mg/L.}

\end{figure}%

\subsection{Bottled water}\label{bottled-water}

Figure~\ref{fig-bottled} and Table~\ref{tbl-bottled} place 142 bottled
waters with a printed bicarbonate or alkalinity on the same scale. They
cover the full range of tap water. Seventeen products have alkalinity
below 15 mg/L and are suitable as diluents; most are purified United
States brands. Thirty-two products have alkalinity between 25 and 60
mg/L and brew like the reference water; because a given product is the
same everywhere, these are suitable as a control water for a validation
study. Forty-two products, many of them European natural mineral waters,
have alkalinity above 150 mg/L. The compilation is not a market census:
supermarket own labels and several large brands publish no analysis, and
45 of the 187 products print neither bicarbonate nor alkalinity.

\begin{longtable}[]{@{}
  >{\raggedright\arraybackslash}p{(\linewidth - 16\tabcolsep) * \real{0.1111}}
  >{\raggedright\arraybackslash}p{(\linewidth - 16\tabcolsep) * \real{0.1111}}
  >{\raggedright\arraybackslash}p{(\linewidth - 16\tabcolsep) * \real{0.1111}}
  >{\raggedright\arraybackslash}p{(\linewidth - 16\tabcolsep) * \real{0.1111}}
  >{\raggedright\arraybackslash}p{(\linewidth - 16\tabcolsep) * \real{0.1111}}
  >{\raggedright\arraybackslash}p{(\linewidth - 16\tabcolsep) * \real{0.1111}}
  >{\raggedright\arraybackslash}p{(\linewidth - 16\tabcolsep) * \real{0.1111}}
  >{\raggedright\arraybackslash}p{(\linewidth - 16\tabcolsep) * \real{0.1111}}
  >{\raggedright\arraybackslash}p{(\linewidth - 16\tabcolsep) * \real{0.1111}}@{}}
\caption{Bottled waters by brewing role, 142 products with bicarbonate
or alkalinity printed on the label.}\label{tbl-bottled}\tabularnewline
\toprule\noalign{}
\begin{minipage}[b]{\linewidth}\raggedright
role (alkalinity, mg/L as CaCO\textsubscript{3})
\end{minipage} & \begin{minipage}[b]{\linewidth}\raggedright
products
\end{minipage} & \begin{minipage}[b]{\linewidth}\raggedright
BR
\end{minipage} & \begin{minipage}[b]{\linewidth}\raggedright
US
\end{minipage} & \begin{minipage}[b]{\linewidth}\raggedright
GB
\end{minipage} & \begin{minipage}[b]{\linewidth}\raggedright
DE
\end{minipage} & \begin{minipage}[b]{\linewidth}\raggedright
FR
\end{minipage} & \begin{minipage}[b]{\linewidth}\raggedright
IT
\end{minipage} & \begin{minipage}[b]{\linewidth}\raggedright
PT
\end{minipage} \\
\midrule\noalign{}
\endfirsthead
\toprule\noalign{}
\begin{minipage}[b]{\linewidth}\raggedright
role (alkalinity, mg/L as CaCO\textsubscript{3})
\end{minipage} & \begin{minipage}[b]{\linewidth}\raggedright
products
\end{minipage} & \begin{minipage}[b]{\linewidth}\raggedright
BR
\end{minipage} & \begin{minipage}[b]{\linewidth}\raggedright
US
\end{minipage} & \begin{minipage}[b]{\linewidth}\raggedright
GB
\end{minipage} & \begin{minipage}[b]{\linewidth}\raggedright
DE
\end{minipage} & \begin{minipage}[b]{\linewidth}\raggedright
FR
\end{minipage} & \begin{minipage}[b]{\linewidth}\raggedright
IT
\end{minipage} & \begin{minipage}[b]{\linewidth}\raggedright
PT
\end{minipage} \\
\midrule\noalign{}
\endhead
\bottomrule\noalign{}
\endlastfoot
diluent (under 15) & 17 & 5 & 5 & 0 & 2 & 0 & 1 & 4 \\
bright (15 to 25) & 12 & 4 & 2 & 1 & 1 & 0 & 1 & 3 \\
reference-like (25 to 60) & 32 & 20 & 8 & 0 & 1 & 1 & 1 & 1 \\
moderate (60 to 150) & 39 & 13 & 7 & 4 & 3 & 6 & 3 & 3 \\
as flat as hard tap water (over 150) & 42 & 4 & 3 & 4 & 17 & 6 & 6 &
2 \\
\end{longtable}

\begin{figure}[H]

\centering{

\pandocbounded{\includegraphics[keepaspectratio]{figures/fig5_bottled.png}}

}

\caption{\label{fig-bottled}Bottled waters by label alkalinity, one row
per country of sale. Shaded regions mark alkalinity below 15 mg/L and
between 25 and 60 mg/L. The products with the lowest and highest
alkalinity in each country are named.}

\end{figure}%

\section{Discussion}\label{discussion}

\subsection{The alkalinity--hardness
relationship}\label{the-alkalinityhardness-relationship}

The single line in Figure~\ref{fig-where} has a geochemical explanation.
Where groundwater or surface water dissolves calcite or dolomite, each
calcium or magnesium ion enters with two bicarbonate ions, so carbonate
hardness and alkalinity are equal when both are expressed as
CaCO\textsubscript{3}. Waters lie above the line when alkalinity comes
with sodium rather than calcium (sodium bicarbonate groundwaters, and
municipal softening, which exchanges calcium for sodium and leaves
bicarbonate in place) and below it when hardness comes from gypsum or
chloride salts rather than carbonate. The fitted slope of log alkalinity
on log hardness is 0.62 rather than 1, which indicates that
non-carbonate hardness makes up a growing share of hardness at the hard
end of the distribution. For brewing this matters in two ways. First,
hardness, which nearly every utility reports, is a usable proxy for
alkalinity, with the factor-of-1.7 uncertainty stated above; the
residual is the price of not measuring. Second, the two variables act on
coffee by different mechanisms, buffering and extraction, and a water
that departs from the line separates them: softened water has the
buffering effect without the cations, and gypsum-rich water the reverse.
These off-line waters are where the two mechanisms could be
distinguished experimentally.

\subsection{The reference target and the
population}\label{the-reference-target-and-the-population}

The Specialty Coffee Association target of about 40 mg/L alkalinity
(Specialty Coffee Association of America, 2009) lies at the upper edge
of the first band, below the population-weighted median of 83 mg/L for
the United States and below the water of about 80 percent of the
population in our data. The target is a laboratory reference, not a
description of tap water, and the model gives one reading of what it
represents: a water that neutralises about 6 percent of a light roast's
acidity, small enough to leave the roast's character intact. Half of the
United States population is between 40 and 80 mg/L, where the
neutralised share is 6 to 12 percent and the predicted sourness change
relative to the reference is under 3 points on a 100-point scale, close
to the resolution of a trained panel. The predicted effect becomes large
only above 150 mg/L, for about 10 percent of the population in our data
and half of the English population served by the three companies with
data. The practical conclusion is that for most people the tap water is
not far from the reference, and for a minority it is far enough that no
recipe adjustment recovers the difference at acceptable cost.

\subsection{Lower sourness is not lower
quality}\label{lower-sourness-is-not-lower-quality}

The model predicts the direction and size of a sourness change. Whether
the change is an improvement depends on who is drinking. In the consumer
study on the same brews, 57 percent of black-coffee drinkers formed a
cluster that preferred low-strength, low-to-medium-extraction coffee and
disliked both bitterness and sourness, while 43 percent formed a cluster
that liked either acidic, fruity coffee or roasted, heavier coffee
(Cotter et al., 2021; Guinard et al., 2023). For the majority cluster, a
water that removes a quarter of a light roast's acidity moves the cup
toward what they prefer; for the minority cluster, which includes most
of the specialty-coffee community that discusses brewing water, it moves
the cup away. The compensation rules in Table~\ref{tbl-rules} are
therefore rules for the second group, and the framing of hard water as a
problem is that group's framing. A brewer serving the first group might
treat band four as favourable. The model does not take sides on this; it
quantifies the shift.

\subsection{Which uncertainty
dominates}\label{which-uncertainty-dominates}

Three uncertainties enter the predictions, and they dominate in
different places. Below about 100 mg/L, the titratable acidity range
(about 10 percent of the central value) dominates: the neutralised share
is small and its uncertainty is the acidity denominator. Between 100 and
250 mg/L, the factor-of-1.8 disagreement between the chemistry and
sensory estimates of compensation dominates the dial-in advice, though
not the neutralised share itself. For Brazil, the imputation uncertainty
dominates everything: a municipality with imputed alkalinity 64 mg/L has
a 68 percent interval of 38 to 109 mg/L, which spans three bands, so
Brazilian band assignments are statements about the distribution of
municipalities, not about any one of them. The 1.8 factor has a
plausible origin: the sourness-per-acid calibration comes from roast
means, and roast changes which acids are present as well as how much.
Citric and malic acid degrade during roasting while quinic acid and
chlorogenic acid lactones form (Batali et al., 2021; Yeager et al.,
2023), and these acids differ in perceived sourness per milliequivalent.
A calibration across strength at fixed roast, which the original data
contain but the published tables do not, would resolve it.

\subsection{Cation effects and dark
roast}\label{cation-effects-and-dark-roast}

In the study of Liu and colleagues, dark roast brewed in the hardest of
three waters was rated more acidic than in the softest (Liu et al.,
2025), the opposite direction to buffering. Their waters differed by 20
mg/L in alkalinity but by 48 mg/L in hardness; acidity was assessed from
relative chromatographic peak areas and an eleven-member recognition
panel rather than titration; and a water supplier was a co-author and
funder. The result is nevertheless consistent with the extraction
mechanism proposed from binding energies (Hendon et al., 2014) and
inconsistent with the null result for organic acids at 100 mg/L
(Bratthäll et al., 2024), and the two experimental studies used
different coffees, roasts and analytical methods. If cation-assisted
extraction is real and roast-dependent, positive for dark roast where
acids are bound in a more degraded matrix and small or negative for
light roast, then the predictions for dark roast in bands four and five
are too pessimistic and the two mechanisms partly cancel for dark roast
in hard water. This is the first hypothesis the validation protocol
should address, and the off-line waters of Section 5.1 are the natural
instrument.

\subsection{Scope}\label{scope}

The model applies to filter brewing at 87 to 93 °C. Brew temperature has
little effect on titratable acidity or on the sensory profile at fixed
strength and extraction (Batali et al., 2020), so the rules do not
depend on it within that range. Espresso is excluded: its extraction
physics, strength and acid balance differ, and the sensory data do not
cover it. Immersion methods extract more slowly and at a fixed ratio;
the acidity data used here are for percolation, and the immersion
studies in the literature give titratable acidity values about 20
percent different at comparable strength, in the direction expected from
their higher strength. The model treats the brew as an ideal solution
and bicarbonate protonation as instantaneous and complete to the extent
the pH allows; it does not model carbon dioxide loss, calcium--acid
complexation, or the pH-dependence of acid extraction from the grounds.
Each of these could shift the neutralised share by a few percentage
points; none is likely to change the band structure.

\subsection{Coverage and data access}\label{coverage-and-data-access}

The atlas covers countries with open data. England is represented by
three companies, two of them in chalk regions; Brazilian alkalinity is
imputed; softened and filtered water in homes is not visible in utility
data. Population-weighted statistics are dominated by the United States
and per-locality statistics by small groundwater systems; both are
reported. Most English alkalinity data exist only in company web lookups
under terms that reserve all rights and were not used; permission from
the companies, or a compilation by a United Kingdom institution holding
database rights, would roughly triple the English coverage. A cheap
improvement for the whole atlas would be for utilities to publish
alkalinity alongside hardness in their annual reports; the measurement
is routine, and its absence is the single largest gap in the data.

\section{Validation protocol}\label{sec-protocol}

The predictions can be tested with domestic equipment. One session
consists of three cups of the same coffee prepared with the same recipe
on the same day: tap water; a bottled water with alkalinity between 25
and 60 mg/L (Table~\ref{tbl-bottled}); and tap water with the dose
increased according to Table~\ref{tbl-rules}. A second person serves the
cups without identification. The taster reports, for each pair, which
cup is more sour, which is more bitter and which is preferred. The
record includes the city and utility (alkalinity taken from the atlas),
the roast level stated on the package, the use of a filter or softener,
and optionally a photograph of an aquarium alkalinity test strip. The
analysis is a mixed-effects logistic regression of the response
``tap-water cup is more sour'' on predicted acid neutralisation, with
roast level as a fixed effect and taster as a random effect; the
comparison of the increased-dose cup with the reference-water cup tests
the compensation rule. With paired forced-choice responses, about 150
sessions distributed across the bands give 80 percent power to detect
the predicted difference of 20 percentage points between the lowest and
highest bands. The protocol, the response form and the analysis code are
in the repository, and the analysis plan will be registered before the
first submission.

\section{Appendix: bottled waters closest to the
reference}\label{sec-appendix}

Table~\ref{tbl-bottled-reference} lists, for each country of sale, the
three bottled waters whose label composition is closest to the reference
water, as a starting point for the control water in Section 6. The full
ranking, including tap-water localities, is in the repository
(docs/rankings.md).

\begin{longtable}[]{@{}
  >{\raggedright\arraybackslash}p{(\linewidth - 10\tabcolsep) * \real{0.1667}}
  >{\raggedright\arraybackslash}p{(\linewidth - 10\tabcolsep) * \real{0.1667}}
  >{\raggedright\arraybackslash}p{(\linewidth - 10\tabcolsep) * \real{0.1667}}
  >{\raggedright\arraybackslash}p{(\linewidth - 10\tabcolsep) * \real{0.1667}}
  >{\raggedright\arraybackslash}p{(\linewidth - 10\tabcolsep) * \real{0.1667}}
  >{\raggedright\arraybackslash}p{(\linewidth - 10\tabcolsep) * \real{0.1667}}@{}}
\caption{Bottled waters closest to the reference water in each country
of sale, by the distance \textbar ln(alkalinity/40)\textbar{} +
0.5·\textbar ln(hardness/68)\textbar; values from
labels.}\label{tbl-bottled-reference}\tabularnewline
\toprule\noalign{}
\begin{minipage}[b]{\linewidth}\raggedright
country
\end{minipage} & \begin{minipage}[b]{\linewidth}\raggedright
product
\end{minipage} & \begin{minipage}[b]{\linewidth}\raggedright
alkalinity
\end{minipage} & \begin{minipage}[b]{\linewidth}\raggedright
hardness
\end{minipage} & \begin{minipage}[b]{\linewidth}\raggedright
pH
\end{minipage} & \begin{minipage}[b]{\linewidth}\raggedright
distance
\end{minipage} \\
\midrule\noalign{}
\endfirsthead
\toprule\noalign{}
\begin{minipage}[b]{\linewidth}\raggedright
country
\end{minipage} & \begin{minipage}[b]{\linewidth}\raggedright
product
\end{minipage} & \begin{minipage}[b]{\linewidth}\raggedright
alkalinity
\end{minipage} & \begin{minipage}[b]{\linewidth}\raggedright
hardness
\end{minipage} & \begin{minipage}[b]{\linewidth}\raggedright
pH
\end{minipage} & \begin{minipage}[b]{\linewidth}\raggedright
distance
\end{minipage} \\
\midrule\noalign{}
\endhead
\bottomrule\noalign{}
\endlastfoot
Brazil & Dias D'Ávila: Água Mineral Dias D'Ávila sem Gás Pet 5L & 42 &
39 & & 0.32 \\
Brazil & Legítima Lindóia: Legítima Lindóia & 39 & 36 & 6.4 & 0.34 \\
Brazil & Cristal (Fonte Santa Bárbara): Água Mineral Cristal & 50 & 53 &
6.6 & 0.35 \\
United States & Crystal Geyser: Crystal Geyser Alpine Spring Water
(bottled at Olancha Spring) & 53 & 59 & 6.8 & 0.36 \\
United States & JUST Water: JUST Water Spring Water & 57 & 60 & 7.8 &
0.41 \\
United States & Crystal Geyser: Crystal Geyser Alpine Spring Water
(bottled at Beaughan Spring) & 46 & 36 & 7.4 & 0.45 \\
United Kingdom & Volvic: Natural Mineral Water 1L & 61 & 63 & 7.0 &
0.46 \\
United Kingdom & FIJI Water: Natural Artesian Water & 104 & 103 & 7.9 &
1.16 \\
United Kingdom & Deeside: Deeside Mineral Water (still) & 21 & 22 & 6.1
& 1.19 \\
Germany & Bad Brambacher: Bad Brambacher Naturell & 76 & 81 & & 0.74 \\
Germany & Black Forest: Black Forest Still & 25 & 27 & & 0.92 \\
Germany & Spreequell: Spreequell Naturell & 11 & 56 & & 1.40 \\
France & Mont Blanc: Source Mont Blanc eau minérale naturelle (6x1.5L) &
53 & 80 & 7.6 & 0.35 \\
France & Volvic: Volvic Nature eau minérale naturelle & 66 & 70 & &
0.51 \\
France & Abatilles: Abatilles eau minérale naturelle plate & 107 & 94 &
8.2 & 1.15 \\
Italy & Levissima: Levissima Naturale (150 cl PET) & 47 & 57 & 7.8 &
0.24 \\
Italy & Conad: Conad acqua minerale naturale oligominerale 1.5 l & 112 &
111 & 7.6 & 1.27 \\
Italy & Dolomia: Acqua Dolomia naturale & 123 & 127 & 8.2 & 1.44 \\
Portugal & Cruzeiro: Água do Cruzeiro & 74 & 70 & 7.2 & 0.62 \\
Portugal & Caramulo: Água Caramulo & 20 & 14 & 6.7 & 1.47 \\
Portugal & Vimeiro: Vimeiro Lisa & 19 & 12 & 6.3 & 1.62 \\
\end{longtable}

\section{Data and code availability}\label{data-and-code-availability}

The atlas, the measurement table with per-row provenance, the
bottled-water table, the model code, the figure and table generators and
this manuscript are in the project repository, under the MIT licence for
code and CC BY 4.0 for produced data. Upstream sources retain their own
licences and are listed with checksums. The redistributable raw
downloads (about 1.5 GB: SISAGUA, TapWaterData, EPA, the English
open-data hub, Northern Ireland Water, IBGE and ONS population files)
are attached to the tagged release, with a manifest per source; files
obtained under terms that do not permit redistribution are not included.

\section{Acknowledgements}\label{acknowledgements}

The author used an AI assistant (Anthropic's Claude) for data
engineering, literature retrieval, modelling and drafting under his
direction; all decisions, interpretations and errors are the author's.
Access to subscription articles was through the CAPES Portal de
Periódicos via the Universidade Federal do Paraná. No funding was
received.

\section*{References}\label{references}
\addcontentsline{toc}{section}{References}

\protect\phantomsection\label{refs}
\begin{CSLReferences}{1}{0}
\bibitem[\citeproctext]{ref-Batali2021}
Batali, M. E., Cotter, A. R., Frost, S. C., Ristenpart, W. D., \&
Guinard, J.-X. (2021). Titratable acidity, perceived sourness, and
liking of acidity in drip brewed coffee. \emph{ACS Food Science \&
Technology}, \emph{1}(4), 559--569.
\url{https://doi.org/10.1021/acsfoodscitech.0c00078}

\bibitem[\citeproctext]{ref-Batali2020temp}
Batali, M. E., Ristenpart, W. D., \& Guinard, J.-X. (2020). Brew
temperature, at fixed brew strength and extraction, has little impact on
the sensory profile of drip brew coffee. \emph{Scientific Reports},
\emph{10}, 16450. \url{https://doi.org/10.1038/s41598-020-73341-4}

\bibitem[\citeproctext]{ref-Bratthall2024}
Bratthäll, T., Figueira, J., \& Nording, M. L. (2024). Influence of
divalent cations on the extraction of organic acids in coffee determined
by GC-MS and NMR. \emph{Heliyon}.

\bibitem[\citeproctext]{ref-Cotter2021}
Cotter, A. R., Batali, M. E., Ristenpart, W. D., \& Guinard, J.-X.
(2021). Consumer preferences for black coffee are spread over a wide
range of brew strengths and extraction yields. \emph{Journal of Food
Science}, \emph{86}(1), 194--205.
\url{https://doi.org/10.1111/1750-3841.15561}

\bibitem[\citeproctext]{ref-Frost2020}
Frost, S. C., Ristenpart, W. D., \& Guinard, J.-X. (2020). Effects of
brew strength, brew yield, and roast on the sensory quality of drip
brewed coffee. \emph{Journal of Food Science}, \emph{85}(8), 2530--2543.
\url{https://doi.org/10.1111/1750-3841.15326}

\bibitem[\citeproctext]{ref-Guinard2023}
Guinard, J.-X., Frost, S., Batali, M., Cotter, A., Lim, L. X., \&
Ristenpart, W. D. (2023). A new coffee brewing control chart relating
sensory properties and consumer liking to brew strength, extraction
yield, and brew ratio. \emph{Journal of Food Science}, \emph{88}(5),
2168--2177. \url{https://doi.org/10.1111/1750-3841.16531}

\bibitem[\citeproctext]{ref-Hendon2014}
Hendon, C. H., Colonna-Dashwood, L., \& Colonna-Dashwood, M. (2014). The
role of dissolved cations in coffee extraction. \emph{Journal of
Agricultural and Food Chemistry}, \emph{62}(21), 4947--4950.
\url{https://doi.org/10.1021/jf501687c}

\bibitem[\citeproctext]{ref-Liu2025}
Liu, Y., Lin, J., Hu, S., Wang, G., Yang, C., Zhang, Z., Ren, D., Yi,
L., \& Li, S. (2025). Exploring the impacts and mechanisms of water on
the taste extractions and perceptions of coffee brews: A case study in
filter brewing. \emph{Journal of Food Composition and Analysis},
\emph{146}. \url{https://doi.org/10.1016/j.jfca.2025.107987}

\bibitem[\citeproctext]{ref-SISAGUA}
Ministério da Saúde do Brasil. (2026). \emph{SISAGUA -- controle
semestral (dados abertos)}.
\url{https://dadosabertos.saude.gov.br/dataset/sisagua-controle-semestral}.

\bibitem[\citeproctext]{ref-SCAwater}
Specialty Coffee Association of America. (2009). \emph{SCAA standard:
Water for brewing specialty coffee}. Specialty Coffee Association of
America.

\bibitem[\citeproctext]{ref-Stream}
Stream. (2026). \emph{Stream water industry open data portal}.
\url{https://www.streamwaterdata.co.uk}.

\bibitem[\citeproctext]{ref-TapWaterData2026}
TapWaterData. (2026). \emph{TapWaterData water hardness dataset}.
\url{https://www.tapwaterdata.com/water-hardness}.

\bibitem[\citeproctext]{ref-EPASYR4}
United States Environmental Protection Agency. (2024). \emph{Six-year
review 4 compliance monitoring data (2012--2019)}.
\url{https://www.epa.gov/dwsixyearreview/six-year-review-4-compliance-monitoring-data-2012-2019}.

\bibitem[\citeproctext]{ref-Yeager2023}
Yeager, S. E., Batali, M. E., Guinard, J.-X., \& Ristenpart, W. D.
(2023). Acids in coffee: A review of sensory measurements and
meta-analysis of chemical composition. \emph{Critical Reviews in Food
Science and Nutrition}, \emph{63}(8), 1010--1036.
\url{https://doi.org/10.1080/10408398.2021.1957767}

\end{CSLReferences}




\end{document}
